Unlike a sales funnel, which is more of a model that focuses on the buyer's journey,
the sales pipeline focuses on the seller's actions and is very process focused.

\medskip

A sales pipeline is important because it promotes the following advantages:

\begin{enumerate}

    % ── Stage 1 ──────────────────────────────────────────────
    \item \textbf{Planning and Research.}
    \begin{enumerate}
        \item I consider this to be something I'm very good at, and something that is
        incredibly important. You need to understand the product, and what problem it
        solves.
        \item You also need to understand the market, similar products, and advantages,
        and trade-offs.
        \item What problem does this product solve, and for whom?
        \item Who are the competitors, and what are the trade-offs between solutions?
    \end{enumerate}

    % ── Stage 2 ──────────────────────────────────────────────
    \item \textbf{Lead Qualification.}
    \begin{enumerate}
        \item Defining criterion for potential prospects (simply, crafting buyer
        personas). Allows for optimized efforts.
    \end{enumerate}

    % ── Stage 3 ──────────────────────────────────────────────
    \item \textbf{Engaging (or Prospecting).}
    \begin{enumerate}
        \item Reach out to people who may want what is being sold.
    \end{enumerate}

    % ── Stage 4 ──────────────────────────────────────────────
    \item \textbf{Discovery.}
    \begin{enumerate}
        \item Revolving around building a human-to-human relationship, asking lots of
        questions, and listening to their needs.
        \item After talking to Artin, having a question-based technique that centers
        around discovering pain points is super important.
    \end{enumerate}

    % ── Stage 5 ──────────────────────────────────────────────
    \item \textbf{Positioning (or Proposal).}
    \begin{enumerate}
        \item Once pain points are found, the product must be positioned in a way that
        can solve these pain points.
        \item How does this product directly address the pain points uncovered?
        \item What makes this solution better than the alternatives they've considered?
    \end{enumerate}

    % ── Stage 6 ──────────────────────────────────────────────
    \item \textbf{Presentation (or Demoing).}
    \begin{enumerate}
        \item Show the product in action.
        \item Does the demo speak to their specific pain points, or is it generic?
        \item What moment in the demo is most likely to create conviction?
    \end{enumerate}

    % ── Stage 7 ──────────────────────────────────────────────
    \item \textbf{Handling Objections.}
    \begin{enumerate}
        \item Address concerns and pushbacks.
        \item Is the objection a dealbreaker or a misunderstanding?
    \end{enumerate}

    % ── Stage 8 ──────────────────────────────────────────────
    \item \textbf{Evaluation.}
    \begin{enumerate}
        \item If the planning stage was done correctly, this should be very
        straightforward. The problem should be very wrought-out, and requirements should
        be met. However, in addition to this, the solution should be presented in a way
        that matches the pain point that was described.
        \item Does the solution meet the requirements defined in Stage 1?
        \item Is the value being communicated in the prospect's own language?
    \end{enumerate}

    % ── Stage 9 ──────────────────────────────────────────────
    \item \textbf{Closing.}
    \begin{enumerate}
        \item What is the decision-making process on their end?
        \item What is the next concrete step?
    \end{enumerate}

    % ── Stage 10 ─────────────────────────────────────────────
    \item \textbf{After-Sales.}
    \begin{enumerate}
        \item What would make them a repeat buyer or an advocate?
    \end{enumerate}

\end{enumerate}